\section{Difusión y brecha del punto de torres}
Trabajamos sobre un soporte regular con punto $p$ de reinas positivo, filas y columnas
de masa uno, $p_e\le C/k$, brecha RC al menos $\gamma>0$ y masas diagonales
$\ell_D,\ell_A\le\epsilon$. Son las condiciones deterministas de C, secciones 2 y 5,
y de B, sección 8; se admite reducir el umbral de $\epsilon$ según $C,\gamma$.
La afirmación se aplica a padres y a hijos que estén en la correspondiente clase exterior.
No afirma que una historia sea regular con una probabilidad dada.

Sea $D$ la incidencia RC con dos entradas positivas por casilla,
$S=D\diag(p)D^t$ y $v=(\mathbf1,-\mathbf1)$. La hipótesis de brecha es sobre $v^\perp$
(o la equivalente para los pesos de B, comparables a $p$ por $1+O(1/k)$).
El proyector sobre las perturbaciones relativas tangentes es
\[
 \Pi=I-D^tS^+D\diag(p).
\]
Es uniformemente acotado en norma del supremo. En efecto, para $\|z\|_\infty\le1$,
$b=D\diag(p)z$ tiene coordenadas acotadas por uno y norma euclídea $O(\sqrt k)$.
La brecha da $\|S^+b\|_2=O(\sqrt k)$. La ecuación por vértice y
$\sum_jp_{ij}^2\le C/k$ dan $\|S^+b\|_\infty=O(1)$, también para columnas.

Módulo sumas de potenciales de fila y columna y constantes, la derivada de la acción
negativa de torres es $\log p+\log(1-p)$; la de reinas es
$\log p+3\log(1-p)-\log h_D-\log h_A$, $h_f=1-\ell_f$.
Al proyectar, su diferencia es $O(\epsilon+1/k)$. Para $P=p(1+z)$, $D\diag(p)z=0$,
la derivada en $z$ de la primera gradiente proyectada es $I+O(\|z\|_\infty+1/k)$
en el espacio tangente. Esto sigue al derivar los dos logaritmos y usar la cota de $\Pi$.
El mapa $z\mapsto z-\Pi\nabla(-H_T)(p(1+z))$ preserva una bola de radio
$C'(\epsilon+1/k)$ y tiene constante de Lipschitz menor que $1/2$ si $\epsilon$ es
suficientemente pequeño y $k$ grande. Su punto fijo da
\[
 P_T/p=1+O(\epsilon+1/k).
\]
La construcción asegura positividad en todo el soporte. Para identificar el maximizador
global basta la concavidad por fila. El coeficiente diagonal de la hessiana negativa
es $(1-2t)/(t(1-t))$. Si una entrada $t_1>1/2$, las demás suman $1-t_1$.
La función $\eta(t)=t(1-t)/(1-2t)$, para $t<1/2$, satisface
$\sum_{j>1}\eta(t_j)\le\eta(1-t_1)=-\eta(t_1)$, ya que $\eta(t)/t$ crece.
Cauchy--Schwarz con la restricción $\sum z_j=0$ prueba la semidefinición positiva
de la hessiana en la fila. Los demás casos se obtienen directamente o por continuidad.
La estricta definitud local en la rama pequeña y la concavidad dan unicidad global.
Por comparación de formas cuadráticas $D\diag(P_T)D^t$ hereda la brecha.
Sus autovalores son $1\pm\sigma_j(P_T)$, de modo que
$\|P_T-\mathbf1\mathbf1^t/k\|_{\rm op}\le 1-\gamma'$ y $P_{T,e}\le C'/k$.

\paragraph{Menores y prefijos de veto.}
Suprimir un número fijo de filas y columnas emparejadas y a lo sumo dos sitios por
cada fila o columna deja en el punto restringido $q_0$ déficits RC $O(1/k)$ y una
brecha positiva. Para verlo, extender por cero un vector de $v^\perp$ y restar la
energía de las aristas perdidas; el coste es $O(\|v'\|_2^2/k)$ para un vector $v'$.
Resolver $D\diag(q_0)D^t y=\mathbf1-Dq_0$ cuesta $O(1/k)$ en supremo por la misma
ecuación por vértice, pues las masas siguen siendo $1+O(1/k)$. El punto
$q_0(1+D^ty)$ es factible y su defecto de gradiente es $O(1/k)$.
La inversión precedente da un cambio relativo $1+O(1/k)$ en cualquier sitio superviviente.
La estimación vale tomando las supresiones directamente desde el padre, sin iterar
un error por cada sitio. Para las marginales verdaderas se usan además las entradas
uniformes de prescripción de las rondas 1, 2 y 9 en esta clase.
